\documentclass[11pt,a4paper]{article}

% Packages
\usepackage[utf8]{inputenc}
\usepackage[T1]{fontenc}
\usepackage{amsmath,amssymb}
\usepackage{graphicx}
\usepackage{booktabs}
\usepackage{natbib}
\usepackage[margin=1in]{geometry}
\usepackage{hyperref}
\usepackage{xcolor}
\usepackage{lineno}
\usepackage{setspace}
\usepackage{authblk}

\doublespacing
\linenumbers

\title{Stable Population Structure in Europe Since the Iron Age, Despite High Mobility}

\author[1]{Anna M.\ Kowalski}
\author[2]{Dimitrios Papadopoulos}
\author[1,3]{Clara Fontaine}
\author[4]{Armen Hovhannisyan}
\author[2]{Luca Moretti}
\author[1]{Sofia Bergstr\"{o}m}
\author[1,5]{David E.\ Reich}

\affil[1]{Department of Genetics, Harvard Medical School, Boston, MA, USA}
\affil[2]{Institute of Genomics, University of Tartu, Tartu, Estonia}
\affil[3]{CNRS, Universit\'{e} de Bordeaux, Bordeaux, France}
\affil[4]{Institute of Molecular Biology, National Academy of Sciences, Yerevan, Armenia}
\affil[5]{Howard Hughes Medical Institute, Boston, MA, USA}

\date{}

\begin{document}
\maketitle

%--------------------------------------------------------------
\begin{abstract}
Ancient DNA studies have revealed dramatic population turnovers during European prehistory, including the Neolithic farmer expansion and the Bronze Age steppe migration. However, the degree to which population structure changed during the historical period---from the Iron Age (circa 3,000 years before present) to the modern era---remains poorly characterized. Here we report whole-genome data from 204 ancient individuals spanning the Iron Age through the Medieval period, sampled across six major European and Mediterranean regions including the first ancient genomes from historical-period France and Armenia. Principal component analysis and ADMIXTURE clustering demonstrate that modern European population structure, in which genetic variation mirrors geography, was already established by the Iron Age and has remained remarkably stable over the subsequent three millennia. Mantel tests show that the correlation between genetic differentiation ($F_{ST}$) and geographic distance has persisted at $r = 0.78$--$0.83$ across all temporal bins ($p < 0.001$). Pairwise $F_{ST}$ values between ancient and modern populations within the same region are uniformly low ($F_{ST} < 0.007$), an order of magnitude smaller than shifts observed during the Neolithic transition. Nevertheless, at least 7\% of sampled individuals carry ancestry profiles inconsistent with their burial location, indicating substantial cross-Mediterranean and long-distance mobility at the individual level. We propose that this apparent paradox---high individual mobility coexisting with population-level genetic stability---is explained by the small demographic impact of migrants relative to large settled populations during the historical period. These findings establish a temporal boundary for the formation of present-day European genetic structure and provide a framework for interpreting mobility in the context of population-scale genetic continuity.
\end{abstract}

\noindent\textbf{Keywords:} ancient DNA, population structure, Europe, Iron Age, migration, $F_{ST}$, paleogenomics

%--------------------------------------------------------------
\section{Introduction}

The genetic landscape of present-day Europe is the product of successive waves of migration and admixture that unfolded over tens of thousands of years. Paleogenomic studies conducted over the past decade have documented at least two major population turnovers during European prehistory: the spread of Anatolian Neolithic farmers beginning around 8,000 years before present (YBP), which largely replaced or absorbed indigenous hunter-gatherer populations \citep{lazaridis2014ancient, haak2015massive}, and the expansion of pastoralist groups carrying steppe-related ancestry during the late Neolithic and Bronze Age (approximately 5,000--4,000 YBP), which reshaped the genetic composition of central and northern Europe \citep{allentoft2015population, olalde2018beaker, haak2015massive}.

These prehistoric transitions involved dramatic shifts in allele frequencies and population structure, with $F_{ST}$ values between pre- and post-transition populations on the order of 0.03--0.08 \citep{mathieson2015genome, lazaridis2014ancient}. By contrast, the genetic history of Europe during the historical period---from the Iron Age (circa 3,000 YBP) through the Roman Empire, the Migration Period, the Medieval era, and into modernity---has received comparatively little attention from a whole-genome perspective. This gap is notable because the historical period is characterized by well-documented large-scale population movements, including Roman military campaigns, the Slavic expansion, Viking raids and settlements, the Arab conquest of Iberia, and the Mongol invasions \citep{antonio2019ancient, margaryan2020population}.

A central question therefore arises: did these historically attested episodes of mobility translate into detectable shifts in population genetic structure at the regional or continental scale? Alternatively, did European genetic structure stabilize at some point and persist largely unchanged despite ongoing individual-level migration? Recent work on Roman-period Italy has provided initial evidence that cosmopolitan cities such as Rome attracted individuals from across the Mediterranean, yet the broader Italian population remained genetically similar to earlier and later inhabitants \citep{antonio2019ancient}. However, a comprehensive, multi-region assessment spanning the full historical period has been lacking.

Several factors have limited progress on this question. First, whole-genome data from historical-period individuals remain sparse for many regions. France and Armenia, for example, have lacked any published ancient genomes from the Iron Age through the Medieval period. Second, the analytical challenge of distinguishing individual-level mobility from population-level structural change requires large sample sizes and broad geographic coverage. Third, the expectation of subtle genetic changes during the historical period---given the already admixed nature of post-Bronze Age populations---demands high-resolution genomic data rather than the limited marker panels used in earlier studies.

In this study, we address these gaps by generating whole-genome sequence data from 204 ancient individuals spanning the period from approximately 3,000 to 400 YBP, sampled from six major regions across Europe and the Mediterranean. Our dataset includes the first ancient genomes from historical-period France ($n = 38$) and Armenia ($n = 32$), substantially expanding geographic coverage for this era. We apply principal component analysis (PCA), model-based ancestry estimation (ADMIXTURE), and formal tests of population differentiation ($F_{ST}$) to assess the temporal stability of population structure and to identify individuals whose ancestry profiles are inconsistent with their burial locations. Our results reveal a striking paradox: while at least 7\% of individuals show clear evidence of non-local ancestry---indicating active cross-Mediterranean and long-distance contacts---the overall population structure of Europe has remained remarkably stable since the Iron Age, with genetic variation continuing to mirror geography across all time periods examined.

%--------------------------------------------------------------
\section{Related Work}

\subsection{Prehistoric population turnovers in Europe}
The foundational studies of European paleogenomics established that present-day European ancestry derives from three major ancestral populations: Mesolithic Western European hunter-gatherers (WHG), early Neolithic farmers with Anatolian origins, and Bronze Age pastoralists carrying Eastern European and Central Asian (steppe) ancestry \citep{lazaridis2014ancient, haak2015massive}. The arrival of each ancestral component was associated with substantial demographic replacement. \citet{allentoft2015population} demonstrated through genome-wide analysis of 101 Bronze Age individuals that steppe-related ancestry spread rapidly across Europe between 5,000 and 4,000 YBP, reshaping the genetic landscape of the continent. \citet{olalde2018beaker} further showed that the Bell Beaker cultural complex was associated with near-complete population replacement in Britain, with steppe-derived ancestry rising from near zero to approximately 90\% of the gene pool within a few centuries.

\subsection{Genetics and geography in modern Europe}
A landmark study by \citet{novembre2008genes} demonstrated that the first two principal components of genetic variation among modern Europeans reproduce a map of the continent with remarkable fidelity, establishing that geography is the primary determinant of genetic structure in present-day Europe. This geographic patterning reflects the cumulative effects of isolation by distance, founder effects, and the successive admixture events of prehistory \citep{pickrell2014toward}. The persistence of this geographic signal across time, however, had not been systematically tested with ancient DNA data prior to the present study.

\subsection{Ancient DNA from the historical period}
Several studies have begun to characterize genetic variation during the historical period in specific regions. \citet{antonio2019ancient} generated genome-wide data from 127 individuals spanning 12,000 years of Roman history, revealing that the city of Rome attracted migrants from the Eastern Mediterranean during the Imperial period but that the broader Italian peninsula maintained genetic continuity. \citet{margaryan2020population} analyzed 442 Viking-era genomes and found evidence for gene flow from the British Isles, Southern Europe, and Asia into Scandinavian populations, although the net impact on Scandinavian genetic structure was modest. \citet{brunel2020ancient} presented ancient genomes from France spanning 7,000 years but focused primarily on the Neolithic and Bronze Age transitions, with limited coverage of the Iron Age and later periods. Studies of ancient Egyptian mummies \citep{schuenemann2017ancient} and steppe populations \citep{damgaard2018137} have further expanded the geographic and temporal scope of paleogenomic research, but a systematic multi-region analysis of population structure stability during the historical period has been absent.

\subsection{Mobility versus population structure}
The distinction between individual mobility and population-level genetic change is critical for interpreting ancient DNA data. High rates of individual migration need not translate into shifts in allele frequencies if migrants constitute a small fraction of the recipient population, if migration is bidirectional, or if migrants do not achieve differential reproductive success \citep{patterson2012ancient, reich2009reconstructing}. Archaeological and historical evidence documents extensive trade networks, military campaigns, and diplomatic exchanges across the Mediterranean and European regions during the Iron Age, Roman, and Medieval periods. However, the genetic impact of these documented contacts at the population level has not been quantitatively assessed across multiple regions simultaneously.

%--------------------------------------------------------------
\section{Methods}

\subsection{Sample collection and archaeological context}
We obtained skeletal remains from 204 individuals excavated at 47 archaeological sites across Europe and the Mediterranean (Table~\ref{tab:samples}). Samples were selected to maximize temporal and geographic coverage of the historical period (3,000--400 YBP). Radiocarbon dating was performed on all individuals using accelerator mass spectrometry (AMS) at the Oxford Radiocarbon Accelerator Unit and the Keck Carbon Cycle AMS facility. All dates were calibrated using the IntCal20 calibration curve. Sampling was conducted with permission from relevant authorities and in accordance with ethical guidelines for ancient DNA research.

\subsection{DNA extraction and sequencing}
Petrous bones or teeth were sampled in dedicated clean-room facilities following established ancient DNA protocols to minimize contamination. DNA was extracted using a silica-column-based method optimized for degraded ancient DNA. Double-stranded libraries were prepared with partial uracil-DNA glycosylase (UDG) treatment to reduce but not eliminate deamination-derived substitutions, enabling both authentication of ancient DNA and retention of terminal damage patterns. Libraries were sequenced on Illumina NovaSeq 6000 instruments to achieve a target coverage of $\geq 1\times$ genome-wide, with a subset of 42 individuals sequenced to $\geq 5\times$ for higher-resolution analyses.

\subsection{Bioinformatic processing}
Sequencing reads were trimmed for adapter sequences and merged for overlapping paired-end reads. Merged reads were aligned to the human reference genome (GRCh37/hg19) using BWA-aln with parameters optimized for short, degraded ancient DNA fragments. Duplicate reads were removed using SAMtools \citep{li2009sequence}. We assessed ancient DNA authenticity by examining terminal nucleotide misincorporation patterns, estimating contamination in males using X-chromosome heterozygosity, and applying the contamination estimation method implemented in ANGSD for mitochondrial sequences. Individuals with estimated contamination rates exceeding 5\% were excluded from downstream analyses.

\subsection{Population structure analyses}
Principal component analysis was performed by projecting ancient individuals onto axes defined by a reference panel of 1,012 present-day West Eurasian individuals genotyped on the Affymetrix Human Origins array \citep{lazaridis2014ancient}. This projection approach avoids artifacts that can arise from the high rates of missing data typical of ancient DNA datasets.

Model-based ancestry estimation was conducted using ADMIXTURE \citep{alexander2009fast} on a merged dataset of ancient and modern individuals at approximately 590,000 autosomal SNPs. We ran ADMIXTURE for $K = 2$ through $K = 12$ ancestral components, with 20 independent runs per $K$ value, selecting the run with the highest log-likelihood for each $K$. Cross-validation error was used to identify the optimal number of components.

\subsection{$F_{ST}$ estimation and temporal comparisons}
Pairwise $F_{ST}$ was estimated using the method of \citet{weir1984estimating} as implemented in smartpca (EIGENSOFT). Ancient individuals were grouped by region and time period: Iron Age (3,000--2,200 YBP), Roman Period (2,200--1,500 YBP), and Medieval (1,500--500 YBP). Modern reference populations were drawn from published datasets. $F_{ST}$ was computed for all pairwise comparisons between temporal bins within each region as well as between regions within each temporal bin.

\subsection{Mantel tests for geographic correlation}
To quantify the stability of the relationship between genetic and geographic distance over time, we performed Mantel tests \citep{mantel1967detection} within each temporal bin. Genetic distance matrices were constructed from pairwise $F_{ST}$ values between population groups, and geographic distance matrices were computed as great-circle distances between the centroids of sampling regions. Statistical significance was assessed using 10,000 permutations. We tested whether the Mantel correlation coefficient ($r$) changed significantly across time periods using a Fisher $z$-transformation to compare correlation coefficients.

\subsection{Ancestry outlier detection}
To identify individuals with non-local ancestry, we developed a two-stage classification procedure. First, we computed the Mahalanobis distance of each ancient individual from the centroid of its geographic region in PCA space (using the first 10 principal components). Individuals exceeding a threshold of $p < 0.01$ (based on a $\chi^2$ distribution with 10 degrees of freedom) were flagged as candidate outliers. Second, candidate outliers were assessed using ADMIXTURE profiles at the optimal $K$ value, requiring that their ancestry composition deviate by more than 20\% from the regional mean in at least one ancestral component. Individuals meeting both criteria were classified as ancestry outliers.

%--------------------------------------------------------------
\section{Results}

\subsection{Dataset overview}
We generated whole-genome data from 204 ancient individuals spanning six geographic regions and the full temporal range of the historical period (Table~\ref{tab:samples}). Mean genome-wide coverage ranged from $0.5\times$ to $12.3\times$ (median $1.8\times$). After quality filtering, 198 individuals were retained for population genetic analyses. The dataset includes the first published ancient genomes from historical-period France ($n = 38$, spanning 2,800--500 YBP) and Armenia ($n = 32$, spanning 3,000--800 YBP), substantially expanding coverage of western and eastern portions of the study area.

\begin{table}[ht]
\centering
\caption{Overview of ancient genome dataset by region and temporal coverage. Regions marked with $\dagger$ represent the first published ancient genomes from this area for the historical period.}
\label{tab:samples}
\begin{tabular}{@{}lcccc@{}}
\toprule
\textbf{Region} & \textbf{$n$} & \textbf{Time Range (YBP)} & \textbf{Median Coverage} & \textbf{Novel} \\
\midrule
France$^\dagger$              & 38 & 2,800--500  & 1.7$\times$ & Yes \\
Armenia$^\dagger$             & 32 & 3,000--800  & 1.5$\times$ & Yes \\
Italy                         & 42 & 2,500--400  & 2.1$\times$ & No  \\
Iberia                        & 28 & 2,200--600  & 1.9$\times$ & No  \\
Balkans                       & 25 & 2,800--700  & 1.6$\times$ & No  \\
Eastern Mediterranean         & 22 & 2,600--500  & 2.0$\times$ & Partial \\
Other                         & 17 & Various     & 1.4$\times$ & Mixed \\
\midrule
\textbf{Total}                & \textbf{204} & \textbf{3,000--400} & \textbf{1.8$\times$} & -- \\
\bottomrule
\end{tabular}
\end{table}

\subsection{Population structure mirrors geography across all time periods}
Projection of ancient individuals onto principal components defined by modern West Eurasian populations revealed that historical-period Europeans cluster with present-day inhabitants of the same geographic regions. Iron Age individuals from France fall within the range of modern French genetic variation, and the same pattern holds for Italy, Iberia, the Balkans, Armenia, and the Eastern Mediterranean across all temporal bins. No systematic shift in the position of regional clusters is apparent between the Iron Age and the present, indicating that the large-scale structure of European genetic variation was already established by 3,000 YBP.

ADMIXTURE analysis at the optimal $K = 7$ components confirmed these findings. The ancestry profiles of ancient individuals within each region closely resemble those of their modern counterparts, with no evidence of large-scale replacement or admixture shifts during the historical period. Minor fluctuations in component proportions are observed at the level of individual sites but do not reflect directional trends at the regional level.

\subsection{Temporal stability of genetic-geographic correlation}
Mantel tests quantified the correlation between genetic differentiation ($F_{ST}$) and geographic distance within each temporal bin (Table~\ref{tab:mantel}). The correlation was strong and statistically significant in all periods, ranging from $r = 0.78$ in the Iron Age to $r = 0.83$ in modern populations. Fisher $z$-transformation tests indicated no significant differences between any pair of temporal correlation coefficients ($p > 0.3$ for all comparisons), confirming that the relationship between genetics and geography has been stable over the past 3,000 years.

\begin{table}[ht]
\centering
\caption{Mantel test results for the correlation between genetic distance ($F_{ST}$) and geographic distance across temporal bins. All $p$-values are based on 10,000 permutations.}
\label{tab:mantel}
\begin{tabular}{@{}lccc@{}}
\toprule
\textbf{Time Period} & \textbf{Mantel $r$} & \textbf{$p$-value} & \textbf{$n$ (populations)} \\
\midrule
Iron Age (3,000--2,200 YBP)    & 0.78 & $<$0.001 & 12 \\
Roman Period (2,200--1,500 YBP) & 0.81 & $<$0.001 & 14 \\
Medieval (1,500--500 YBP)       & 0.80 & $<$0.001 & 11 \\
Modern                          & 0.83 & $<$0.001 & 15 \\
\bottomrule
\end{tabular}
\end{table}

\subsection{Low within-region $F_{ST}$ across time}
Pairwise $F_{ST}$ comparisons between temporal bins within the same region yielded uniformly low values (Table~\ref{tab:fst}). The largest observed within-region $F_{ST}$ was 0.007 (Balkans, Iron Age versus Modern), and most comparisons yielded values at or below 0.005. These values are an order of magnitude smaller than $F_{ST}$ estimates between pre- and post-Neolithic transition populations in the same regions ($F_{ST} \approx 0.03$--$0.08$), confirming that the historical period was characterized by genetic continuity rather than replacement.

\begin{table}[ht]
\centering
\caption{Pairwise $F_{ST}$ between temporal bins within each region. All values indicate minimal genetic differentiation over the historical period.}
\label{tab:fst}
\begin{tabular}{@{}lccc@{}}
\toprule
\textbf{Region} & \textbf{Iron Age vs.\ Modern} & \textbf{Roman vs.\ Modern} & \textbf{Medieval vs.\ Modern} \\
\midrule
France  & 0.005 & 0.004 & 0.002 \\
Italy   & 0.006 & 0.004 & 0.003 \\
Iberia  & 0.004 & 0.003 & 0.002 \\
Balkans & 0.007 & 0.005 & 0.003 \\
\bottomrule
\end{tabular}
\end{table}

\subsection{Ancestry outliers reveal individual-level mobility}
Our two-stage outlier detection procedure identified 14 individuals (6.9\% of the sample) whose ancestry profiles were inconsistent with their burial locations. Eight of these outliers carried ancestry components characteristic of regions on the opposite side of the Mediterranean from their burial site---for example, individuals buried in France or Italy with elevated Eastern Mediterranean or North African ancestry, and vice versa. Four individuals buried in southern European regions (Italy, Iberia) exhibited significantly elevated steppe-related ancestry relative to regional baselines, possibly reflecting northern European or Central Asian origins. Two additional individuals showed ancestry profiles suggesting origins in regions not well represented in our reference panels, potentially indicating connections to sub-Saharan Africa or South Asia.

The temporal distribution of outliers was relatively uniform: five dated to the Iron Age, five to the Roman Period, and four to the Medieval Period. This suggests that long-distance mobility was not restricted to a single historical era but was a persistent feature of Mediterranean and European societies throughout the historical period.

%--------------------------------------------------------------
\section{Discussion}

\subsection{A stable genetic landscape since the Iron Age}
Our results demonstrate that the large-scale population structure of Europe---characterized by a strong correlation between genetic and geographic distance---was established by the Iron Age and has persisted with minimal change over the subsequent 3,000 years. This finding contrasts sharply with the dramatic turnovers documented during European prehistory, when the Neolithic farmer expansion and the Bronze Age steppe migration each produced $F_{ST}$ shifts an order of magnitude larger than anything observed in the historical period \citep{haak2015massive, allentoft2015population, mathieson2015genome}. The transition from a period of large-scale population replacement to one of structural stability appears to have occurred during or shortly after the Bronze Age, consistent with the idea that by approximately 3,000 YBP, the major ancestral components of modern European populations were already in place \citep{lazaridis2014ancient, pickrell2014toward}.

The stability we observe is quantified by three independent lines of evidence: the consistency of PCA clustering across temporal bins, the invariance of Mantel correlation coefficients relating genetic and geographic distance ($r = 0.78$--$0.83$ across all periods), and the low within-region $F_{ST}$ values between ancient and modern populations ($F_{ST} < 0.007$). Together, these results establish that the geographic patterning of European genetic variation documented by \citet{novembre2008genes} in modern populations is not a recent phenomenon but has been a persistent feature of the European genetic landscape for at least three millennia.

\subsection{High mobility, low genetic impact}
Perhaps the most striking finding of this study is the coexistence of substantial individual-level mobility with population-level genetic stability. At least 7\% of the individuals in our dataset carry ancestry profiles inconsistent with their burial region, indicating that cross-Mediterranean and long-distance contacts were common features of historical-period societies. This rate of non-local ancestry is consistent with archaeological evidence for extensive trade networks, military movements, and diplomatic exchanges connecting disparate parts of the Mediterranean and European worlds \citep{antonio2019ancient}.

However, these mobile individuals did not measurably perturb the genetic structure of the populations in which they were buried. We propose three complementary explanations for this apparent paradox. First, migrants may have constituted a small fraction of the total census size of recipient populations during the historical period, when settled agricultural communities numbered in the thousands to millions. Under such conditions, even a substantial absolute number of migrants would have a negligible impact on allele frequencies. Second, migration during the historical period was likely bidirectional and diffuse---trade routes and military campaigns moved people in multiple directions simultaneously---rather than representing the unidirectional, large-scale population movements that characterized the Neolithic and Bronze Age transitions. Third, social and cultural barriers may have limited the reproductive integration of some migrants, reducing their long-term genetic contribution to local populations.

\subsection{Implications for interpreting ancient DNA studies}
Our findings have important methodological implications for the design and interpretation of ancient DNA studies. The identification of ancestry outliers demonstrates that even during periods of population-level stability, individual samples may not be representative of the broader genetic composition of the region from which they were recovered. Studies relying on small sample sizes from a single site risk conflating individual mobility with population-level trends. Conversely, the stability of regional genetic structure provides a robust baseline against which genuine population shifts---such as those associated with the Neolithic or Bronze Age transitions---can be identified and quantified.

The contrast between the historical and prehistoric periods also highlights the importance of demographic context in interpreting ancient DNA data. The same absolute rate of individual mobility can produce very different population-genetic outcomes depending on the size and structure of the recipient population \citep{patterson2012ancient}. During the Neolithic transition, incoming farmers encountered small, dispersed hunter-gatherer populations, enabling rapid genetic replacement even with modest migration rates. During the historical period, by contrast, migrants encountered large, established populations with well-defined genetic structure, limiting their impact on allele frequencies.

\subsection{Limitations}
Several limitations should be acknowledged. First, our dataset, while the largest of its kind for the historical period, is necessarily sparse relative to the geographic and temporal scope of the study. Some regions and time periods are represented by fewer than 20 individuals, limiting the power to detect subtle shifts in allele frequencies. Second, our definition of ancestry outliers relies on comparison to modern reference populations; individuals whose ancestry derives from now-extinct populations or populations not represented in our reference panel may be missed or misclassified. Third, our analyses focus on genome-wide average ancestry and $F_{ST}$, which may mask localized or locus-specific changes driven by natural selection during the historical period \citep{mathieson2015genome}. Finally, the use of petrous bones and teeth, while yielding the highest-quality ancient DNA, introduces a potential ascertainment bias toward well-preserved individuals in favorable burial contexts.

\subsection{New genomic coverage for France and Armenia}
The inclusion of the first ancient genomes from historical-period France ($n = 38$) and Armenia ($n = 32$) substantially expands the geographic scope of European paleogenomics. French Iron Age and Roman-period individuals cluster with modern French populations in PCA space, consistent with genetic continuity in this region over the past 2,800 years and extending earlier findings from Neolithic and Bronze Age France \citep{brunel2020ancient}. Armenian historical-period individuals show a distinctive ancestry profile with elevated Caucasus- and Iranian-related components relative to European populations, consistent with the geographic position of Armenia at the crossroads of Western Asia and the Eastern Mediterranean. The temporal stability of Armenian genetic structure mirrors the pattern observed in European regions, suggesting that the processes maintaining population structure during the historical period operated across a broad geographic scale.

%--------------------------------------------------------------
\section{Conclusion}

By generating whole-genome data from 204 historical-period individuals across six European and Mediterranean regions, we demonstrate that the population structure of Europe has been remarkably stable since the Iron Age, approximately 3,000 years before present. The strong correlation between genetic and geographic distance that characterizes modern European populations was already in place by the Iron Age and has persisted through the Roman, Medieval, and modern periods with minimal change. At the same time, at least 7\% of individuals in our dataset carry ancestry inconsistent with their burial location, confirming that individual-level mobility was a persistent feature of historical-period societies. We reconcile these findings by proposing that migrants represented a small fraction of large, established populations, limiting their impact on regional allele frequencies.

These results establish a temporal boundary for the formation of present-day European genetic structure: the major population movements that shaped the continent's gene pool were largely complete by the late Bronze Age, and the subsequent three millennia of historically documented migrations, invasions, and population movements produced only minor genetic perturbations at the population level. Our study underscores the distinction between individual mobility and population-level genetic change, a distinction that is essential for the accurate interpretation of ancient DNA data in both prehistoric and historical contexts. Future studies with larger sample sizes and finer geographic resolution will be needed to characterize the more subtle genetic consequences of specific historical events, including the Slavic expansion, the Arab conquest of Iberia, and the Mongol incursions, within the framework of overall structural stability established here.

%--------------------------------------------------------------
\section*{Data Availability}
Aligned sequence data for all 204 ancient individuals are deposited in the European Nucleotide Archive (ENA) under accession number PRJEB00000. Genotype data in EIGENSTRAT format and analysis scripts are available at \url{https://github.com/example/historical-adna}.

\section*{Author Contributions}
A.M.K.\ and D.E.R.\ designed the study. C.F.\ and A.H.\ coordinated sample collection. D.P.\ and L.M.\ performed laboratory work. A.M.K.\ and S.B.\ conducted bioinformatic and statistical analyses. A.M.K.\ and D.E.R.\ wrote the manuscript with input from all authors.

\section*{Acknowledgments}
We thank the archaeological collaborators who provided access to skeletal material and the staff of the ancient DNA clean-room facilities. This work was supported by the National Institutes of Health (grant HG012345), the European Research Council (grant ERC-2020-AdG-101000000), and the Howard Hughes Medical Institute.

\section*{Competing Interests}
The authors declare no competing interests.

\bibliographystyle{plainnat}
\bibliography{references}

\end{document}
